% preamble.tex -- shared style for the vSLAM course
% Colour scheme: blue = definitions, red = theorems, amber = propositions
% and lemmas, green = "In words" intuition, grey = practice / code.

\usepackage[T1]{fontenc}
\usepackage[utf8]{inputenc}
\usepackage{lmodern}
\usepackage{microtype}
\usepackage[a4paper,margin=2.5cm]{geometry}
\usepackage{amsmath,amssymb,mathtools,bm}
\usepackage{graphicx}
\usepackage{xcolor}
\usepackage{tikz}
\usetikzlibrary{arrows.meta,positioning,calc,shapes.geometric,fit,
  decorations.pathreplacing,3d}
\usepackage{enumitem}
\usepackage{booktabs}
\usepackage[section]{placeins}
\usepackage{listings}
\usepackage[most]{tcolorbox}
\usepackage[colorlinks=true,linkcolor=blue!50!black,
  citecolor=green!40!black,urlcolor=blue!60!black]{hyperref}
\usepackage[nameinlink,capitalise]{cleveref}

% ---------------------------------------------------------------- colours
\definecolor{defblue}{RGB}{30,90,170}
\definecolor{thmred}{RGB}{175,35,35}
\definecolor{propamber}{RGB}{200,130,0}
\definecolor{wordsgreen}{RGB}{40,130,70}
\definecolor{codegrey}{RGB}{90,90,90}
\definecolor{codebg}{RGB}{248,248,246}
\definecolor{lawngreen}{RGB}{70,140,40}

% ----------------------------------------------------- theorem-like boxes
\tcbset{
  basestyle/.style={enhanced,breakable,sharp corners=south,
    fonttitle=\bfseries,boxrule=0.6pt,left=6pt,right=6pt,
    top=4pt,bottom=4pt,separator sign={.}},
}
\newtcbtheorem[number within=chapter]{definition}{Definition}{
  basestyle,label type=definition,
  colback=defblue!5,colframe=defblue,coltitle=white}{def}
\newtcbtheorem[use counter from=definition]{theorem}{Theorem}{
  basestyle,label type=theorem,
  colback=thmred!5,colframe=thmred,coltitle=white}{thm}
\newtcbtheorem[use counter from=definition]{proposition}{Proposition}{
  basestyle,label type=proposition,
  colback=propamber!7,colframe=propamber,coltitle=white}{prop}
\newtcbtheorem[use counter from=definition]{lemma}{Lemma}{
  basestyle,label type=lemma,
  colback=propamber!7,colframe=propamber,coltitle=white}{lem}
\newtcbtheorem[use counter from=definition]{example}{Example}{
  basestyle,label type=example,
  colback=white,colframe=black!40,coltitle=white}{ex}

\newtcolorbox{inwords}{enhanced,breakable,colback=wordsgreen!6,
  colframe=wordsgreen,boxrule=0.6pt,sharp corners=south,
  title={\textbf{In words}},fonttitle=\bfseries,coltitle=white}
\newtcolorbox{practice}[1][]{enhanced,breakable,colback=codebg,
  colframe=codegrey,boxrule=0.6pt,sharp corners=south,
  title={\textbf{Practice}\ifx\relax#1\relax\else\ --- #1\fi},
  fonttitle=\bfseries,coltitle=white}
\newtcolorbox{mower}{enhanced,breakable,colback=lawngreen!6,
  colframe=lawngreen!70!black,boxrule=0.6pt,sharp corners=south,
  title={\textbf{On the lawn}},fonttitle=\bfseries,coltitle=white}

\crefname{definition}{Definition}{Definitions}
\crefname{theorem}{Theorem}{Theorems}
\crefname{proposition}{Proposition}{Propositions}
\crefname{lemma}{Lemma}{Lemmas}
\crefname{example}{Example}{Examples}

\newenvironment{proof}[1][Proof]{\par\noindent\textit{#1.}\ }
  {\hfill$\square$\par\medskip}

% --------------------------------------------------------------- listings
\lstdefinestyle{py}{language=Python,basicstyle=\ttfamily\small,
  backgroundcolor=\color{codebg},frame=single,rulecolor=\color{codegrey},
  keywordstyle=\color{defblue}\bfseries,commentstyle=\color{wordsgreen},
  stringstyle=\color{thmred},showstringspaces=false,breaklines=true,
  columns=fullflexible,keepspaces=true}
\lstset{style=py}

% ---------------------------------------------------------------- macros
\newcommand{\R}{\mathbb{R}}
\newcommand{\N}{\mathbb{N}}
\newcommand{\SO}[1]{\mathrm{SO}(#1)}
\newcommand{\SE}[1]{\mathrm{SE}(#1)}
\newcommand{\so}[1]{\mathfrak{so}(#1)}
\newcommand{\se}[1]{\mathfrak{se}(#1)}
\newcommand{\GL}[1]{\mathrm{GL}(#1)}
\newcommand{\hatop}[1]{#1^{\wedge}}
\newcommand{\veeop}[1]{#1^{\vee}}
\newcommand{\Exp}{\operatorname{Exp}}
\newcommand{\Log}{\operatorname{Log}}
\newcommand{\Ad}{\operatorname{Ad}}
\newcommand{\ad}{\operatorname{ad}}
\newcommand{\tr}{\operatorname{tr}}
\newcommand{\diag}{\operatorname{diag}}
\newcommand{\T}{^{\mathsf{T}}}
\newcommand{\E}{\mathbb{E}}
\newcommand{\Normal}{\mathcal{N}}
\newcommand{\vect}[1]{\bm{#1}}
\newcommand{\x}{\vect{x}}
\newcommand{\z}{\vect{z}}
\newcommand{\uu}{\vect{u}}
\newcommand{\m}{\vect{m}}
\newcommand{\p}{\vect{p}}
\newcommand{\btheta}{\vect{\theta}}
\newcommand{\bphi}{\vect{\phi}}
\newcommand{\brho}{\vect{\rho}}
\newcommand{\bxi}{\vect{\xi}}
\newcommand{\bomega}{\vect{\omega}}
\newcommand{\norm}[1]{\left\lVert #1\right\rVert}
\newcommand{\code}[1]{\texttt{#1}}
